\documentclass[reqno]{amsart}
\usepackage[T1]{fontenc}
\usepackage{amsmath,amssymb,amsthm}
\usepackage{xr-hyper}
\usepackage[hidelinks]{hyperref}
\externaldocument[da-]{20261126T030300--the-data__workbook}
\externaldocument[ob-]{20261126T040100--spaces-forms-and-twists__workbook}
\externaldocument[bp-]{20261126T020210--bessel-poincare-series-and-their-gram-operators__workbook}
\externaldocument[ss-]{20261126T030400--the-seesaw-identity__workbook}
\externaldocument[d6-]{20261126T050300--the-size-of-the-composite-lift__workbook}
\externaldocument[ta-]{20261126T030600--theorem-a-short-forms__workbook}
\externaldocument[a1-]{20261126T080100--free-levels-and-short-bases__workbook}

% Chapter number; \ref{chap} gives it.
\newcounter{chap}
% Chapter numbers, generated from 20261125T012000--grh-workbook__guide.org
% by scripts/make-chapters.py.  Do not edit: change the order of the
% index and rerun the script (./preview-site.sh and CI do this).
%
% Usage in a sheet, after \newcounter{chap}:
%   % Chapter numbers, generated from 20261125T012000--grh-workbook__guide.org
% by scripts/make-chapters.py.  Do not edit: change the order of the
% index and rerun the script (./preview-site.sh and CI do this).
%
% Usage in a sheet, after \newcounter{chap}:
%   % Chapter numbers, generated from 20261125T012000--grh-workbook__guide.org
% by scripts/make-chapters.py.  Do not edit: change the order of the
% index and rerun the script (./preview-site.sh and CI do this).
%
% Usage in a sheet, after \newcounter{chap}:
%   \input{chapters}\chapnumber{<sheet ID>}
\makeatletter
\providecommand\chapters@set[2]{\@namedef{chapters@#1}{#2}}
\chapters@set{20261125T001000}{1}% zeros-to-ramp-sums
\chapters@set{20261125T001100}{2}% inert-primes-and-estimate-b
\chapters@set{20261125T005000}{3}% galois-fourier-and-averaging
\chapters@set{20261125T005500}{4}% columns-from-a-sphere
\chapters@set{20261125T000000}{5}% slope-method
\chapters@set{20261126T030100}{6}% a-twist-for-each-prime
\chapters@set{20261125T003000}{7}% the-chain-and-the-square-root
\chapters@set{20261125T010000}{8}% rank-lemma-and-zero-estimate
\chapters@set{20261125T010500}{9}% poincare-series-and-near-orthogonality
\chapters@set{20261124T223000}{10}% elliptic-curves-for-the-proof
\chapters@set{20261125T009500}{11}% fock-model-and-theta-intertwiners
\chapters@set{20261124T223100}{12}% complex-multiplication-and-inert-reduction
\chapters@set{20261125T008000}{13}% one-cm-determinant
\chapters@set{20261124T214500}{14}% supersingular-theta-divisibility
\chapters@set{20261124T214600}{15}% the-isotropic-plane
\chapters@set{20261125T013700}{16}% pairings-on-a-formal-group-of-height-two
\chapters@set{20261125T013500}{17}% hecke-nondegeneracy-on-the-tree
\chapters@set{20261126T030200}{18}% the-bound-is-attained
\chapters@set{20261125T009000}{19}% kronecker-limit-and-chowla-selberg
\chapters@set{20261125T002000}{20}% definite-theta-lifts
\chapters@set{20261125T015500}{21}% siegel-weil-in-the-smallest-case
\chapters@set{20261125T002100}{22}% definite-theta-lifts-and-norms
\chapters@set{20261125T016500}{23}% hecke-operators-and-theta-lifts
\chapters@set{20261125T016700}{24}% divisibility-along-the-chain
\chapters@set{20261125T004000}{25}% denominators-and-sha
\chapters@set{20261125T004100}{26}% tate-twists-and-sha
\chapters@set{20261126T020430}{27}% the-minus-part-by-gauss-sums
\chapters@set{20261125T014100}{28}% congruence-modules
\chapters@set{20261125T014500}{29}% eisenstein-congruences-and-selmer-classes
\chapters@set{20261125T014700}{30}% pseudocharacters-and-the-length-bound
\chapters@set{20261125T011000}{31}% galois-representations-black-box
\chapters@set{20261126T030300}{32}% the-data
\chapters@set{20261126T020210}{33}% bessel-poincare-series-and-their-gram-operators
\chapters@set{20261126T030400}{34}% the-seesaw-identity
\chapters@set{20261126T030500}{35}% divisibility-of-the-composite-lift
\chapters@set{20261126T030600}{36}% theorem-a-short-forms
\chapters@set{20261126T030700}{37}% the-determinant-of-a-chain
\chapters@set{20261126T030800}{38}% the-cost-of-one-lift
\chapters@set{20261126T030900}{39}% denominators-of-projections-to-lifts
\chapters@set{20261126T031000}{40}% twists-with-few-congruences
\chapters@set{20261126T031100}{41}% theorem-c-and-the-main-inequality
\providecommand\chapnumber[1]{%
  \@ifundefined{chapters@#1}%
    {\PackageError{chapters}{Sheet #1 is not listed in the index}%
      {List it in 20261125T012000--grh-workbook__guide.org\MessageBreak
       and run scripts/make-chapters.py.}}%
    {\edef\thechap{\@nameuse{chapters@#1}}}}
\makeatother
\chapnumber{<sheet ID>}
\makeatletter
\providecommand\chapters@set[2]{\@namedef{chapters@#1}{#2}}
\chapters@set{20261125T001000}{1}% zeros-to-ramp-sums
\chapters@set{20261125T001100}{2}% inert-primes-and-estimate-b
\chapters@set{20261125T005000}{3}% galois-fourier-and-averaging
\chapters@set{20261125T005500}{4}% columns-from-a-sphere
\chapters@set{20261125T000000}{5}% slope-method
\chapters@set{20261126T030100}{6}% a-twist-for-each-prime
\chapters@set{20261125T003000}{7}% the-chain-and-the-square-root
\chapters@set{20261125T010000}{8}% rank-lemma-and-zero-estimate
\chapters@set{20261125T010500}{9}% poincare-series-and-near-orthogonality
\chapters@set{20261124T223000}{10}% elliptic-curves-for-the-proof
\chapters@set{20261125T009500}{11}% fock-model-and-theta-intertwiners
\chapters@set{20261124T223100}{12}% complex-multiplication-and-inert-reduction
\chapters@set{20261125T008000}{13}% one-cm-determinant
\chapters@set{20261124T214500}{14}% supersingular-theta-divisibility
\chapters@set{20261124T214600}{15}% the-isotropic-plane
\chapters@set{20261125T013700}{16}% pairings-on-a-formal-group-of-height-two
\chapters@set{20261125T013500}{17}% hecke-nondegeneracy-on-the-tree
\chapters@set{20261126T030200}{18}% the-bound-is-attained
\chapters@set{20261125T009000}{19}% kronecker-limit-and-chowla-selberg
\chapters@set{20261125T002000}{20}% definite-theta-lifts
\chapters@set{20261125T015500}{21}% siegel-weil-in-the-smallest-case
\chapters@set{20261125T002100}{22}% definite-theta-lifts-and-norms
\chapters@set{20261125T016500}{23}% hecke-operators-and-theta-lifts
\chapters@set{20261125T016700}{24}% divisibility-along-the-chain
\chapters@set{20261125T004000}{25}% denominators-and-sha
\chapters@set{20261125T004100}{26}% tate-twists-and-sha
\chapters@set{20261126T020430}{27}% the-minus-part-by-gauss-sums
\chapters@set{20261125T014100}{28}% congruence-modules
\chapters@set{20261125T014500}{29}% eisenstein-congruences-and-selmer-classes
\chapters@set{20261125T014700}{30}% pseudocharacters-and-the-length-bound
\chapters@set{20261125T011000}{31}% galois-representations-black-box
\chapters@set{20261126T030300}{32}% the-data
\chapters@set{20261126T020210}{33}% bessel-poincare-series-and-their-gram-operators
\chapters@set{20261126T030400}{34}% the-seesaw-identity
\chapters@set{20261126T030500}{35}% divisibility-of-the-composite-lift
\chapters@set{20261126T030600}{36}% theorem-a-short-forms
\chapters@set{20261126T030700}{37}% the-determinant-of-a-chain
\chapters@set{20261126T030800}{38}% the-cost-of-one-lift
\chapters@set{20261126T030900}{39}% denominators-of-projections-to-lifts
\chapters@set{20261126T031000}{40}% twists-with-few-congruences
\chapters@set{20261126T031100}{41}% theorem-c-and-the-main-inequality
\providecommand\chapnumber[1]{%
  \@ifundefined{chapters@#1}%
    {\PackageError{chapters}{Sheet #1 is not listed in the index}%
      {List it in 20261125T012000--grh-workbook__guide.org\MessageBreak
       and run scripts/make-chapters.py.}}%
    {\edef\thechap{\@nameuse{chapters@#1}}}}
\makeatother
\chapnumber{<sheet ID>}
\makeatletter
\providecommand\chapters@set[2]{\@namedef{chapters@#1}{#2}}
\chapters@set{20261125T001000}{1}% zeros-to-ramp-sums
\chapters@set{20261125T001100}{2}% inert-primes-and-estimate-b
\chapters@set{20261125T005000}{3}% galois-fourier-and-averaging
\chapters@set{20261125T005500}{4}% columns-from-a-sphere
\chapters@set{20261125T000000}{5}% slope-method
\chapters@set{20261126T030100}{6}% a-twist-for-each-prime
\chapters@set{20261125T003000}{7}% the-chain-and-the-square-root
\chapters@set{20261125T010000}{8}% rank-lemma-and-zero-estimate
\chapters@set{20261125T010500}{9}% poincare-series-and-near-orthogonality
\chapters@set{20261124T223000}{10}% elliptic-curves-for-the-proof
\chapters@set{20261125T009500}{11}% fock-model-and-theta-intertwiners
\chapters@set{20261124T223100}{12}% complex-multiplication-and-inert-reduction
\chapters@set{20261125T008000}{13}% one-cm-determinant
\chapters@set{20261124T214500}{14}% supersingular-theta-divisibility
\chapters@set{20261124T214600}{15}% the-isotropic-plane
\chapters@set{20261125T013700}{16}% pairings-on-a-formal-group-of-height-two
\chapters@set{20261125T013500}{17}% hecke-nondegeneracy-on-the-tree
\chapters@set{20261126T030200}{18}% the-bound-is-attained
\chapters@set{20261125T009000}{19}% kronecker-limit-and-chowla-selberg
\chapters@set{20261125T002000}{20}% definite-theta-lifts
\chapters@set{20261125T015500}{21}% siegel-weil-in-the-smallest-case
\chapters@set{20261125T002100}{22}% definite-theta-lifts-and-norms
\chapters@set{20261125T016500}{23}% hecke-operators-and-theta-lifts
\chapters@set{20261125T016700}{24}% divisibility-along-the-chain
\chapters@set{20261125T004000}{25}% denominators-and-sha
\chapters@set{20261125T004100}{26}% tate-twists-and-sha
\chapters@set{20261126T020430}{27}% the-minus-part-by-gauss-sums
\chapters@set{20261125T014100}{28}% congruence-modules
\chapters@set{20261125T014500}{29}% eisenstein-congruences-and-selmer-classes
\chapters@set{20261125T014700}{30}% pseudocharacters-and-the-length-bound
\chapters@set{20261125T011000}{31}% galois-representations-black-box
\chapters@set{20261126T030300}{32}% the-data
\chapters@set{20261126T020210}{33}% bessel-poincare-series-and-their-gram-operators
\chapters@set{20261126T030400}{34}% the-seesaw-identity
\chapters@set{20261126T030500}{35}% divisibility-of-the-composite-lift
\chapters@set{20261126T030600}{36}% theorem-a-short-forms
\chapters@set{20261126T030700}{37}% the-determinant-of-a-chain
\chapters@set{20261126T030800}{38}% the-cost-of-one-lift
\chapters@set{20261126T030900}{39}% denominators-of-projections-to-lifts
\chapters@set{20261126T031000}{40}% twists-with-few-congruences
\chapters@set{20261126T031100}{41}% theorem-c-and-the-main-inequality
\providecommand\chapnumber[1]{%
  \@ifundefined{chapters@#1}%
    {\PackageError{chapters}{Sheet #1 is not listed in the index}%
      {List it in 20261125T012000--grh-workbook__guide.org\MessageBreak
       and run scripts/make-chapters.py.}}%
    {\edef\thechap{\@nameuse{chapters@#1}}}}
\makeatother
\chapnumber{20261126T080200}
\renewcommand{\thesection}{\thechap.\arabic{section}}

\theoremstyle{definition}
\newtheorem{problem}{Problem}[section]
\theoremstyle{plain}
\newtheorem{theorem}{Theorem}[section]
\theoremstyle{definition}
\newtheorem{definition}[theorem]{Definition}

\newcommand{\Z}{\mathbb{Z}}
\newcommand{\Q}{\mathbb{Q}}
\newcommand{\R}{\mathbb{R}}
\newcommand{\C}{\mathbb{C}}
\newcommand{\F}{\mathbb{F}}
\newcommand{\Gb}{\bar G}
\newcommand{\Sram}{S_{\mathrm{ram}}}
\newcommand{\Aut}{\operatorname{Aut}}
\newcommand{\diag}{\operatorname{diag}}
\newcommand{\ps}{\mathsf p_s}
\newcommand{\Av}{\mathrm{Av}}
\newcommand{\Nc}{\mathcal N}
\newcommand{\tl}[1]{\texttt{#1}}

\setlength{\emergencystretch}{2em}

\title{The good part and the ramified forms}
\author{}
\date{November 26, 2026}

\begin{document}
\maketitle
\refstepcounter{chap}\label{chap}

\paragraph{Prerequisites.}
Characters of finite groups: orthogonality, the central idempotents of the group algebra and integrality of character values \cite{Se}.  From Chapter~\ref{da-chap}: the data (Definition~\ref{da-def:data}) and Problems~\ref{da-pr:degree}, \ref{da-pr:T}, \ref{da-pr:m-big} and~\ref{da-pr:MN}.  From Chapter~\ref{ob-chap}: the chosen places \(\mathcal M\), the group \(\Gb\), the idempotent \(e_0\) and the order \(\Nc\) (Section~\ref{ob-sec:chosen}); the action \(\rho\) of \(\Gb\) on the principal forms \(H(Y_0)\) and the target \(T = e_0H(Y_0)\) (Section~\ref{ob-sec:forms}); the gauge (Definition~\ref{ob-def:gauge}); the source forms \(S\), \(D_X = \dim S\), the operators \(T_w\), the scalars \(z_w\) and \(\Sram\) (Section~\ref{ob-sec:source}); tame characters and the lower chain \(\Theta_\xi\) (Sections~\ref{ob-sec:twists} and~\ref{ob-sec:chain}, Definition~\ref{ob-def:final-level}).  From Chapter~\ref{ss-chap}: the source level \(K^X_0\) and Problem~\ref{ss-pr:ram-image}.  From Chapter~\ref{ta-chap}: the norm \(\lVert\cdot\rVert_\sigma\) and the inputs listed in Section~\ref{ta-sec:theorem-a}.  From Chapter~\ref{d6-chap}: the class count.  From Chapter~\ref{a1-chap}: free levels and their forms.

\medskip
Theorem~\ref{ta-thm:A} of Chapter~\ref{ta-chap} uses two inputs that this chapter proves.  The normalized translates of Theorem~A apply \(\#\Gb\,\rho(e_0)\rho(s)\) to principal forms.  Why does the factor \(\#\Gb\) suffice for integrality, and how much can the operator enlarge norms?  And the ramified source forms are cut out by \(f\) conditions, one at each prime \(w \mid \pi_N\): how large is \(\Sram\)?  The answers are input~3 (Theorem~\ref{thm:good-part}) and input~5 (Theorem~\ref{thm:ramified-dim}).  For the first we use the compatibility of the lower chain with translations, read from the proof, and the class count of Chapter~\ref{d6-chap}; for the second, free levels.

\section{Central idempotents}\label{sec:idempotents}

In this section \(G\) is a finite group, and an element \(e = \sum_u e(u)\,u\) of \(\C[G]\) has coefficients \(e(u)\).  The central idempotents are the sums of the idempotents \((\dim\mathsf V/\#G)\sum_u \operatorname{tr}(u^{-1}\mid\mathsf V)\,u\) of distinct irreducible representations \(\mathsf V\) \cite{Se}.  Let \(S_3\) act on \(\C^3\) by permuting coordinates, let \(\C^3_0 \subset \C^3\) be the plane of vectors with coordinate sum zero, and let \(e\) be the idempotent of \(\C^3_0\).  It is a central idempotent, not the \(e_0\) of any data.

\begin{problem}\label{pr:s3}
\(6e\) has coefficient \(4\) at the identity, \(0\) at the transpositions and \(-2\) at the \(3\)-cycles.  One has \(e(1, 0, 0) = (\tfrac23, -\tfrac13, -\tfrac13)\) and \(6e(1, 0, 0) = (4, -2, -2)\), of squared length \(24\).  The largest \(3\)-adic absolute value of a coefficient of \(e\) is \(3\); that of \(6e\) is \(1\).
\end{problem}

\begin{problem}\label{pr:character}
For a central idempotent \(e\) of \(\C[G]\), the function \(u \mapsto \#G\,e(u)\) is the character of a representation of \(G\) of dimension at most \(\#G\).  In Problem~\ref{pr:s3} it is the character of \(\C^3_0 \oplus \C^3_0\), of dimension \(4\).
\end{problem}

\begin{problem}\label{pr:coeff-bounds}
Let \(e\) be a central idempotent of \(\C[G]\) and \(u \in G\).  Then \(\lvert\eta(\#G\,e(u))\rvert_p \le 1\) for every prime \(p\) and every ring homomorphism \(\eta\colon\C\to\C_p\), and \(\lvert\sigma(e(u))\rvert \le 1\) for every \(\sigma \in \Aut(\C)\).
\end{problem}

\section{The good translate}\label{sec:good}

\begin{definition}\label{def:good-translate}
For \(s \in \Gb\) the \emph{good translate} is the linear map \(\ps := \#\Gb\,\rho(e_0)\rho(s)\colon H(Y_0) \to T\).
\end{definition}

By Definition~\ref{ta-def:normalized}, the normalized translate satisfies \(A_s\varphi = c_B\,c_{\mathrm{ss}}^{-1}\,\ps(\mathcal P\varphi)\) for every \(\varphi \in \Sram\).  The proof relates the lower chain of \(\ps z\) to that of \(z\) through the action of \(\Gb\).  For a tame character \(\xi\) and \(s \in \Gb\) it fixes an element \(\tilde s_\xi\) of \(K^\circ(Y_m)\) that reduces to \(\diag(s_w, 1)\) at each chosen place \(w\), where \(s = (s_w)_{w \in \mathcal M}\), and to \(1\) at every \(w \in Q_\xi\), and shows that for all \(y \in H(Y_0)\) and \(g \in G_f(Y_m)\)
\[
  \Theta_\xi(\rho(s)y)(g) = (\Theta_\xi y)(g\,\tilde s_\xi) .
\]
\emph{Status}: read from the proof \cite{GL}.  In this chapter \(u\) denotes an element of \(\Gb\), and for every element of \(\Gb\), such as \(us\), the element fixed in this way is written with a tilde and the subscript \(\xi\).

\begin{problem}\label{pr:expansion}
For every tame character \(\xi\), \(z \in H(Y_0)\), \(s \in \Gb\) and \(g \in G_f(Y_m)\),
\[
  (\Theta_\xi\,\ps z)(g) = \sum_{u \in \Gb} \#\Gb\,e_0(u)\,(\Theta_\xi z)(g\,\widetilde{us}_\xi) .
\]
\end{problem}

Integrality is measured by the gauge.  We use the following classical fact.  \emph{Integral Weyl modules.}  The Weyl module of a highest weight has a basis, dual to the products of minors indexed by semistandard tableaux, in which an invertible matrix \(h\) with entries in a commutative ring \(A\) and \(\det h \in A^\times\) acts by a matrix with entries in \(A\) \cite{DEP}, \cite{Fu}.  \emph{Status}: classical.  The proof takes the coordinates of Definition~\ref{ob-def:gauge} dual to a chosen integral basis of products of minors, in which the action of such matrices is again integral \cite{GL}.

\begin{problem}\label{pr:gauge}
Fix a prime \(p\), a ring homomorphism \(\eta\colon \C \to \C_p\), a framed space \(V\) and a point \(g\) of its finite adelic group.  For Weyl vectors \(\omega, \omega'\) and \(\lambda \in \C\): \(\lVert\lambda\omega\rVert_g = \lvert\eta(\lambda)\rvert_p\lVert\omega\rVert_g\), \(\lVert\omega + \omega'\rVert_g \le \max(\lVert\omega\rVert_g, \lVert\omega'\rVert_g)\), and \(\lVert\omega\rVert_{gh} = \lVert\omega\rVert_g\) for every \(h \in K^\circ(V)\).
\end{problem}

\begin{problem}\label{pr:good-integral}
Let \(p\) be a prime, \(\eta\colon\C\to\C_p\) a ring homomorphism, \(\xi\) a tame character, \(\lambda \in \C\), \(z \in H(Y_0)\) and \(s \in \Gb\).  If \(\lambda\,\Theta_\xi z\) is integral at \(\eta\), so is \(\lambda\,\Theta_\xi\,\ps z\).
\end{problem}

In Problem~\ref{pr:s3} the denominator \(3\) of \(e\) is cleared by \(\#G = 6\): this is the role of the factor \(\#\Gb\).  The price is paid in size.

\begin{problem}\label{pr:pointwise}
Let \(z \in H(Y_0)\), \(s \in \Gb\), \(\sigma \in \Aut(\C)\) and \(c \in \R\).  If \(\lvert\sigma(\Theta_1 z)(g)\rvert^2 \le e^c\) for every \(g \in G_f(Y_m)\), then \(\lvert\sigma(\Theta_1\,\ps z)(g)\rvert^2 \le (\#\Gb)^4e^c\) for every \(g\).  In the model of Problem~\ref{pr:s3}, with the identity in place of \(\Theta_1\), \(z = (1, 0, 0)\) and \(c = 0\), for every \(s \in S_3\) the coordinates of \(6e\,s(1, 0, 0)\) have absolute value at most \(\sum_u \lvert 6e(u)\rvert = 8 \le (\#G)^2 = 36\), and the largest is \(4\).
\end{problem}

Recall from Chapter~\ref{d6-chap} the class count, read from the proof: for the data, let \(G_f\) be the finite adelic unitary group of the standard definite hermitian space of rank at most \(n_m\) and \(\Gamma\) its rational points; if a level \(K\) contains the principal congruence subgroup of \(\pi_M\pi_N\), then \(\#(\Gamma\backslash G_f/K) \le e^{x^{30}}\).

\begin{problem}\label{pr:count}
The final level \(K_{\mathrm{last}}(1)\) of Definition~\ref{ob-def:final-level} contains the principal congruence subgroup of \(\pi_M\pi_N\) in \(K^\circ(Y_m)\), and \(Y_m\) has rank at most \(n_m\).  Hence, for \(y \in H(Y_0)\), \(\sigma \in \Aut(\C)\) and \(c \in \R\): if \(\lvert\sigma(\Theta_1 y)(g)\rvert^2 \le e^c\) for all \(g \in G_f(Y_m)\), then \(\lVert\Theta_1 y\rVert_\sigma^2 \le e^{x^{30} + c}\).
\end{problem}

\begin{problem}\label{pr:group-size}
\(\#\Gb\) divides \(\Nc\), and \(x^{30} + 4\log\Nc \le x^{99}\).
\end{problem}

Problems~\ref{pr:pointwise}--\ref{pr:group-size} give the first estimate and Problems~\ref{pr:expansion}--\ref{pr:good-integral} the second.  This is input~3 of Section~\ref{ta-sec:theorem-a}, word for word (\(\#\Gb\,\rho(e_0)\rho(s)z = \ps z\)).

\begin{theorem}[the good part]\label{thm:good-part}
Let \(z \in H(Y_0)\) and \(s \in \Gb\).  If \(\lvert\sigma(\Theta_1 z)(g)\rvert^2 \le e^c\) for all \(g \in G_f(Y_m)\), with \(c \in \R\) and a fixed \(\sigma \in \Aut(\C)\), then
\[
  \lVert\Theta_1(\#\Gb\,\rho(e_0)\rho(s)z)\rVert_\sigma^2 \le e^{x^{99} + c}.
\]
For every prime \(p\), every ring homomorphism \(\eta\colon \C \to \C_p\), every tame character \(\xi\) and every \(\lambda \in \C\), if \(\lambda\,\Theta_\xi z\) is integral at \(\eta\), so is \(\lambda\,\Theta_\xi(\#\Gb\,\rho(e_0)\rho(s)z)\).
\end{theorem}

\section{The ramified forms}\label{sec:ramified}

For \(w \mid \pi_N\) put \(\Av_w := \sum_{i=0}^{N-2}T_w^i\) (in this chapter).  Then \(\Sram = \bigcap_w \ker\Av_w\), and by Problem~\ref{ss-pr:ram-image}, \(\Sram\) is the image of \(\prod_{w \mid \pi_N}(1 - T_w)\).  We first count dimensions in two finite models, in which the source is replaced by the functions on a finite set and \(T_w\) by translations.

\begin{problem}\label{pr:fp-instance}
Let \(F_1\) be the operator \((F_1v)(y) = v(2y)\) on complex-valued functions \(v\) on \(\F_5^\times\), so \(F_1^4 = 1\).  Then \(\ker(1 + F_1 + F_1^2 + F_1^3)\) is the space of vectors of coordinate sum zero, it equals \(\operatorname{im}(1 - F_1)\), and its dimension is \(3 = (1 - \tfrac14)\cdot 4\); the fixed vectors of \(F_1\) are the constants.  Moreover, for \(F\) in a ring with identity and an integer \(n \ge 0\), \(n - \sum_{i<n}F^i = (1 - F)\sum_{i<n}\sum_{j<i}F^j\); for \(n = 4\) the last factor is \(3 + 2F + F^2\).
\end{problem}

\begin{problem}\label{pr:z4-instance}
On complex-valued functions on \((\Z/4)^2\) let \(F_1, F_2\) be the translations by \((1, 0)\) and \((0, 1)\).  Then \(\bigcap_k\ker(1 + F_k + F_k^2 + F_k^3)\) has dimension \(9\), \(\ker(1 - F_1)\) has dimension \(4\), and \(\ker\bigl((1 - F_1)(1 - F_2)\bigr)\) has dimension \(7 \le 4 + 4\).  With two operators in place of the \(f\) places and order \(4\) in place of \(N - 1\), the bound \((1 - f/(N - 1))D_X\) of Theorem~\ref{thm:ramified-dim} becomes \((1 - \tfrac24)\cdot16 = 8 \le 9\).
\end{problem}

\begin{problem}\label{pr:kernel-sum}
For endomorphisms \(F, F'\) of a finite-dimensional vector space, \(\dim\ker FF' \le \dim\ker F + \dim\ker F'\).
\end{problem}

Write \(K^X_0\langle z_w\rangle\) for the subgroup generated by \(K^X_0\) and the scalar matrix \(z_w\).

\begin{problem}\label{pr:order}
Let \(w \mid \pi_N\) and \(i \in \Z_{\ge 0}\).  The scalar \(z_w^i\) lies in \(K^X_0\) if and only if \(N - 1\) divides \(i\).  Hence \(T_w^{N-1} = 1\) on \(S\), and \([K^X_0\langle z_w\rangle : K^X_0] = N - 1\).  Formally, for \(N = 5\) (no data have \(N = 5\)) and \(z_w\) reducing to \(2\) at \(w\), the least \(i \ge 1\) with \(z_w^i \in K^X_0\) is \(4\).
\end{problem}

For the standard definite hermitian spaces over \(E\), recall from Chapter~\ref{a1-chap} that a level \(K\) is \emph{free} if \(\Gamma \cap gKg^{-1} = 1\) for every finite adelic \(g\), \(\Gamma\) the rational points; that, for the data, every level of rank at most \(4n_m\) contained in the integral matrices that reduce to the identity at every \(w \in \mathcal M\) is free (Problem~\ref{a1-pr:free}); and that for free levels \(K \subseteq K'\) the forms of level \(K\) of a given weight have \([K' : K]\) times the dimension of those of level \(K'\) (Problem~\ref{a1-pr:dim-free}).

\begin{problem}\label{pr:fixed-dim}
For \(w \mid \pi_N\), the source forms fixed by \(T_w\) are exactly the forms of the weight of \(S\) and level \(K^X_0\langle z_w\rangle\); this level and \(K^X_0\) are free; and \(\dim\ker(1 - T_w) = D_X/(N - 1)\).
\end{problem}

\begin{problem}\label{pr:ram-dim}
\(\dim\ker\prod_{w \mid \pi_N}(1 - T_w) \le fD_X/(N - 1)\), and \(\dim \Sram \ge (1 - f/(N - 1))D_X\).
\end{problem}

\begin{problem}\label{pr:f-N}
For data, \(f \le b\) and \(N - 1 \ge b^9\); hence \(f/(N - 1) \le b^{-8}\).
\end{problem}

Problems~\ref{pr:order}--\ref{pr:f-N}, with Problem~\ref{pr:kernel-sum}, prove input~5 of Section~\ref{ta-sec:theorem-a}.

\begin{theorem}[the ramified forms]\label{thm:ramified-dim}
\(\dim \Sram \ge (1 - f/(N - 1))D_X\), and \(f/(N - 1) \le b^{-8}\).
\end{theorem}

In the proof of Theorem~\ref{bp-thm:B}, the chosen finite set \(J \subseteq \Gb\) satisfies \(\#J\,D_X \le 5d_T\log(1 + 2b^2)\), where \(d_T = \dim_\C T\); for this \(J\), Theorem~\ref{thm:ramified-dim} and Problem~\ref{bp-pr:loss} give \(\#J\,(D_X - \dim\Sram) \le 2b^{-2}d_T\).

We have proved Theorem~A's inputs~3 and~5.  Taken as given: the compatibility of the lower chain with the action of \(\Gb\) (read from the proof), the integral Weyl modules (classical), the class count of Chapter~\ref{d6-chap} (read from the proof), and the freeness of the levels and the dimension of forms at free levels (Chapter~\ref{a1-chap}).  Input~1 of Theorem~A, that the \(e^U\)-short forms in \(\Sram\) span \(\Sram\), is deduced in Problem~\ref{ta-pr:ramify-short} from the statement that the \(4^{-f}e^U\)-short forms in \(S\) span \(S\), which Chapter~\ref{a1-chap} proves (Theorem~\ref{a1-thm:short-span}).

\paragraph{In the Lean formalization.}
{\raggedright In the Lean proof \cite{GL} (prefix \tl{New.Data.} omitted): Definition~\ref{def:good-translate} is \tl{goodProjTranslate}; the statement read in Section~\ref{sec:good} is \tl{lowerChain\_principalRep}; the coordinates used in the gauge are indexed by \tl{New.WeylIndexFamily}.  For the entries of the Weyl action matrix, see \tl{New.weylRep\_entry\_mem}.  The class count recalled from Chapter~\ref{d6-chap} is the level case of \tl{card\_stdClasses\_le\_exp}, which allows any subgroup containing the specified principal congruence subgroup.  Theorem~\ref{thm:good-part} is \tl{normSqAt\_goodProjTranslate\_le} with \tl{goodProjTranslate\_mem\_integralForms}.  The proof bounds the fixed space only from above, \(\dim\ker(1 - T_w) \le D_X/(N - 1)\) (\tl{finrank\_ker\_one\_sub\_scalingGenerator\_le}); the equalities of Problems~\ref{pr:order} and~\ref{pr:fixed-dim} are derived here.  Theorem~\ref{thm:ramified-dim} is \tl{le\_finrank\_ramifiedSourceForms} with \tl{f\_div\_N\_sub\_one\_le}.\par}

\section*{Hints}

\subsection*{\S\ref{sec:idempotents}}
\begin{description}
\item[\ref{pr:s3}] The character of \(\C^3_0\) is \(2, 0, -1\) on the three classes, and \(e\) acts on \(\C^3 = \C(1,1,1) \oplus \C^3_0\) as the projection onto \(\C^3_0\).
\item[\ref{pr:character}] Write \(e\) as a sum of idempotents of irreducible representations; which representation has character \(\dim\mathsf V\cdot\operatorname{tr}(u^{-1}\mid\mathsf V)\)?  Count dimensions with \(\sum_{\mathsf V}(\dim\mathsf V)^2 = \#G\).
\item[\ref{pr:coeff-bounds}] Character values are sums of roots of unity, and \(\lvert\operatorname{tr}(u\mid\mathsf V)\rvert \le \dim\mathsf V\).  Applying \(\sigma\) to the coefficients of \(e\) gives a central idempotent.
\end{description}

\subsection*{\S\ref{sec:good}}
\begin{description}
\item[\ref{pr:expansion}] Expand \(\rho(e_0)\rho(s) = \sum_u e_0(u)\rho(us)\) and use linearity of \(\Theta_\xi\).
\item[\ref{pr:gauge}] \((gh)^{-1} = h^{-1}g^{-1}\), and \(h\), \(h^{-1}\) are integral at every finite place; use the integral Weyl modules.
\item[\ref{pr:good-integral}] Problems~\ref{pr:expansion}, \ref{pr:coeff-bounds} and~\ref{pr:gauge} at the points \(g\) and \(g\,\widetilde{us}_\xi\).
\item[\ref{pr:pointwise}] The triangle inequality for the Weyl norm, and Problem~\ref{pr:coeff-bounds} for \(\sigma(e_0(u))\).
\item[\ref{pr:count}] At \(w \in \mathcal M\), \(w\) divides \(\pi_M\); \(Q_1 = \varnothing\).  Each class has weight at most \(1\).
\item[\ref{pr:group-size}] \(g \mapsto \diag(g, 1)\) embeds \(\mathrm{GL}_r(k_w)\) in \(\mathrm{GL}_{rt}(k_w)\).  Problem~\ref{ob-pr:group-order}, and \(x \ge 2\).
\end{description}

\subsection*{\S\ref{sec:ramified}}
\begin{description}
\item[\ref{pr:fp-instance}] Fixed vectors are the functions on the orbits.  Multiplication by \(2\) permutes \(\F_5^\times\) in one cycle.  For the identity, \(\sum_{j<i}F^j\) times \(1 - F\) telescopes.
\item[\ref{pr:z4-instance}] Functions on \((\Z/4)^2\) are \(\C^4 \otimes \C^4\), and \(F_1, F_2\) act on one factor each; compute \(\operatorname{rank}\bigl((1 - F_1)(1 - F_2)\bigr)\).  Which functions are fixed by both translations?
\item[\ref{pr:kernel-sum}] Apply rank and nullity to the restriction of \(F'\) to \(\ker(FF')\).
\item[\ref{pr:order}] The chosen \(z_w\) is an integral unit and equals \(1\) off \(w, \bar w\).  Which last-row congruence for \(z_w^i\) remains at the places dividing \(\pi_N\)?  Scalars are central, so \(K^X_0\langle z_w\rangle\) is the union of the cosets \(z_w^iK^X_0\).
\item[\ref{pr:fixed-dim}] A form fixed by \(K^X_0\) and by \(z_w\) is fixed by the group they generate; \(z_w\) is \(1\) at \(\mathcal M\), and \(X_0\) has rank \(k \le n_m\).  Then Problem~\ref{pr:order}.
\item[\ref{pr:ram-dim}] Problems~\ref{pr:kernel-sum} and~\ref{pr:fixed-dim}, the count of the \(w\) in Problem~\ref{ta-pr:ramify-short}, and rank and nullity for \(\prod_w(1 - T_w)\).
\item[\ref{pr:f-N}] \(f = T(s)\varphi(q_0)\), \(\varphi(q_0) \le q_0 \le C\), and \(x^{10^{12}} \le b\); for \(N\), Problems~\ref{da-pr:MN} and~\ref{da-pr:m-big}.
\end{description}

\begin{thebibliography}{DEP}
\bibitem[DEP]{DEP} C.~De~Concini, D.~Eisenbud and C.~Procesi, \emph{Young diagrams and determinantal varieties}, Invent.\ Math.\ \textbf{56} (1980), 129--165.
\bibitem[Fu]{Fu} W.~Fulton, \emph{Young Tableaux}, London Math.\ Soc.\ Student Texts 35, Cambridge Univ.\ Press, 1997.
\bibitem[GL]{GL} L.~Sernau, \emph{GRH for Dirichlet \(L\)-functions, in Lean~4}, a shorter rewriting of the Lean proof, text by Claude, \url{https://github.com/lsernau/grh-lean}, commit \texttt{72fead1}, readable edition \texttt{grh/}.
\bibitem[Se]{Se} J.-P.~Serre, \emph{Linear Representations of Finite Groups}, Grad.\ Texts in Math.\ 42, Springer, 1977.
\end{thebibliography}

\end{document}
